\section{Implementation}
\label{sec:impl}
ShadowLedger is implemented in Go as a set of small, independently tested packages: erasure
coding (over an RS library), rendezvous placement, the Proof-of-Storage engine, the chain
manager (fork-choice tree, reorg, finality, state), the uint64 VM and the SHL compiler, the
registration-PoW and faucet-PoW primitives, and the peer-to-peer layer. The network is deployed
as a live prototype with a bootstrap node on a public cloud host plus additional peers;
discovery uses DNS seeds and peer exchange with no central coordinator. The deterministic core
is exercised by an automated test suite spanning roughly sixteen packages. Table~\ref{tab:modules}
lists the principal modules and their responsibilities; the separation keeps each mechanism
independently testable, which is how the security and functional checks of \S\ref{sec:eval} are
obtained in isolation before integration.

\begin{table}[t]
\centering
\caption{Principal implementation modules and responsibilities.}
\label{tab:modules}
\small
\begin{tabular}{ll}
\toprule
module & responsibility \\
\midrule
erasure      & RS encode/decode, shard hashing, reconstruction \\
placement    & HRW shard assignment, holder selection \\
consensus    & Proof-of-Storage leader election, round fallback \\
chain        & fork-choice tree, reorg, finality, block apply \\
state        & accounts, validator registry, contract storage \\
vm + shl     & uint64 VM, gas metering, SHL compiler \\
regpow       & one-shot registration PoW (Sybil gate) \\
faucet       & onboarding PoW, anchor verification \\
p2p          & DNS-seed/peer-exchange discovery, gossip, sync \\
\bottomrule
\end{tabular}
\end{table}

\textbf{Networking and discovery.} Membership is maintained without any central directory. A new
node bootstraps from a small set of DNS seeds baked into the binary, then learns further peers by
exchange (each peer shares a sample of its own peer set), so the view of membership grows and heals
under churn, which is the same set the placement module feeds to HRW. Two channels carry traffic: a
control channel gossips headers and full blocks for state application, and a shard-transfer channel
serves and requests individual shards for reconstruction. Genesis is derived deterministically from
the embedded parameters rather than gossiped, so every node that ships the same binary computes an
identical block~$0$ and therefore an identical chain of prev-hash links, which is what lets a
zero-configuration node join a running network safely.

\textbf{Deployment notes.} The prototype runs as zero-configuration mainnet: network parameters
(genesis, seed nodes, erasure policy) are baked into the binary, so a freshly installed node joins
automatically, in the spirit of how reference clients ship their genesis and seeds. The bootstrap
node runs on a single small cloud instance behind a standard firewall; the only operational
prerequisite is opening the control/transfer ports. Software updates are \emph{advisory}: a node
reports a newer release but never auto-applies it, since silent auto-update would be a central
kill switch in a system whose point is decentralization. We treat the live network as a testnet:
the tokens carry no real value, which lets us exercise consensus, reorg, finality, the faucet, and
contract deployment on a genuinely public deployment without inviting financial risk before an
external audit.
